\begin{theorem}[Watatani 指标与屏幕补全成本的函子独立性]\label{thm:conclusion-watatani-screen-completion-functor-independence}
\leanverified{paper\_conclusion\_watatani\_screen\_completion\_functor\_independence}
沿用命题 \ref{prop:conclusion-watatani-handle-identity-trace-moment} 的记号，记
\[
\mathrm{Ind}_{\mathrm W}(E_m)=\sum_{x\in X_m}d_m(x)\,\widetilde P_x.
\]
则
\[
\tau_{m,1}\!\bigl(\mathrm{Ind}_{\mathrm W}(E_m)\bigr)
=
\frac{1}{2^m}\sum_{x\in X_m}d_m(x)^2
=
\frac{S_2(m)}{2^m}.
\]
然而对同一折叠系统，屏幕补全成本
\[
\operatorname{Cost}(S)=c(\Gamma_S)-1
\]
却可随屏幕 \(S\) 从 \(0\) 变化到 \(2^{m(n-1)}\)。因此，不存在仅依赖于 \(\mathrm{Ind}_{\mathrm W}(E_m)\) 的函子，能够恢复所有屏幕的补全成本。
\end{theorem}
\begin{proof}
第一式正是命题 \ref{prop:conclusion-watatani-handle-identity-trace-moment} 的迹矩化公式在 \(g=0\) 的直接改写；因此 \(\mathrm{Ind}_{\mathrm W}(E_m)\) 与其归一化迹只依赖于纤维多重度多重集 \(\{d_m(x)\}_{x\in X_m}\)，与屏幕选择无关。

另一方面，对同一折叠系统，若取全屏 \(S=\overrightarrow{\mathcal F}^{\,n-1}_m\)，则由定理 \ref{thm:conclusion-fullscreen-partialscreen-complexity-transition} 有
\[
\operatorname{Cost}(S)=0,
\]
而轴向屏幕由推论 \ref{cor:conclusion-axial-screen-deficit-exponential-law} 给出
\[
\operatorname{Cost}(S^{(1)})=2^{m(n-1)}.
\]
二者对应同一 \(\mathrm{Ind}_{\mathrm W}(E_m)\)，却具有完全不同的补全成本。故任何只经由 \(\mathrm{Ind}_{\mathrm W}(E_m)\) 因子化的函子，都不可能恢复全部屏幕的 \(\operatorname{Cost}(S)\)。证毕。
\end{proof}
